% =========================================================
% Proof: Additivity of Infima
% Source: volume-iii/analysis/bounding/notes/notes-bounds-arithmetic.tex
% =========================================================

\subsection*{Additivity of Infima}
\label{prf:additivity-of-infimums}

\begin{remark*}[Return]
$\leftarrow$ \hyperref[thm:inf-add]{Back to Theorem in Notes}
\end{remark*}

\begin{theorem*}[Additivity of Infimum]
Let \(A\) and \(B\) be nonempty subsets of \(\mathbb{R}\) that are bounded below. Define
\[
A+B := \{a+b : a\in A,\ b\in B\}.
\]
Then
\[
\inf(A+B)=\inf A+\inf B.
\]
\end{theorem*}

\begin{proof}
Let
\[
\alpha:=\inf A,\qquad \beta:=\inf B.
\]

We first show that \(\alpha+\beta\) is a lower bound for \(A+B\).
If \(x\in A+B\), then \(x=a+b\) for some \(a\in A\) and \(b\in B\).
Since \(\alpha=\inf A\), we have \(\alpha\le a\) for all \(a\in A\).
Since \(\beta=\inf B\), we have \(\beta\le b\) for all \(b\in B\).
Adding these inequalities gives
\[
\alpha+\beta\le a+b=x.
\]
Thus \(\alpha+\beta\) is a lower bound for \(A+B\). Therefore,
\[
\alpha+\beta\le \inf(A+B).
\]

Now let \(\varepsilon>0\) be arbitrary.
Because \(\alpha=\inf A\), there exists \(a_\varepsilon\in A\) such that
\[
a_\varepsilon<\alpha+\frac{\varepsilon}{2}.
\]
Because \(\beta=\inf B\), there exists \(b_\varepsilon\in B\) such that
\[
b_\varepsilon<\beta+\frac{\varepsilon}{2}.
\]
Then \(a_\varepsilon+b_\varepsilon\in A+B\), so by the definition of infimum,
\[
\inf(A+B)\le a_\varepsilon+b_\varepsilon.
\]
Using the inequalities above,
\[
a_\varepsilon+b_\varepsilon
<
\left(\alpha+\frac{\varepsilon}{2}\right)
+
\left(\beta+\frac{\varepsilon}{2}\right)
=
\alpha+\beta+\varepsilon.
\]
Hence
\[
\inf(A+B)<\alpha+\beta+\varepsilon.
\]
Since \(\varepsilon>0\) was arbitrary, it follows that
\[
\inf(A+B)\le \alpha+\beta.
\]

Combining this with the earlier inequality
\[
\alpha+\beta\le \inf(A+B),
\]
we conclude that
\[
\inf(A+B)=\alpha+\beta.
\]
That is,
\[
\inf(A+B)=\inf A+\inf B.
\]
\end{proof}

\begin{remark*}[Proof shape]
Two-part argument: exhibit the candidate value as an upper (lower) bound,
then show no smaller (larger) value qualifies using the
$\varepsilon$-characterization of the supremum (infimum).
\end{remark*}

\begin{remark*}[Dependencies]
Definition of supremum (infimum); \hyperref[prop:eps-char]{$\varepsilon$-Characterization, Proposition~\ref*{prop:eps-char}}.
\end{remark*}
